\subsection{Promeasures on a locally convex space}

Let E be a locally convex space. We denote by \(\mathcal{F}(\mathrm{E})\) the set of closed linear subspaces of E of finite codimension, ordered by the relation \(\supset\) . For every \(V \in \mathcal{F}(\mathrm{E})\) , \(p_{V}\) denotes the canonical mapping of E onto E/V. Let V and W be two elements of \(\mathcal{F}(\mathrm{E})\) such that \(V \supset W\) ; we denote by \(p_{VW}\) the mapping of E/W into E/V deduced from the identity mapping of E by passage to the quotients. The family \(\mathcal{Q}(\mathrm{E}) = (\mathrm{E}/\mathrm{V}, p_{\mathrm{VW}})\) is an inverse system of locally convex spaces, indexed by \(\mathcal{F}(\mathrm{E})\) . It is called the inverse system of finite-dimensional quotients of E.

It can be shown that the inverse limit of the inverse system \(\mathcal{Q}(\mathbf{E})\) is canonically isomorphic to the algebraic dual \(E^{\prime*}\) of \(E^{\prime}\) , equipped with the weak topology \(\sigma(\mathbf{E}^{\prime*},\mathbf{E}^{\prime})\) .

DEFINITION 1. --- Let E be a locally convex space. One calls promeasure on E every inverse system \(^{(1)}\) of measures (§4, No.~2, Def. 1) on the inverse system of finite-dimensional quotients of E.

In other words, a promeasure \(\mu\) on E is a family \((\mu_{\mathrm{V}})_{\mathrm{V}\in \mathcal{F}(\mathrm{E})}\) , where \(\mu_{\mathrm{V}}\) is a bounded (positive) measure on the finite-dimensional space \(\mathrm{E / V}\) , and where \(\mu_{\mathrm{V}} = p_{\mathrm{VW}}(\mu_{\mathrm{W}})\) when \(\mathrm{V}\supset \mathrm{W}\) . All of the measures \(\mu_{\mathrm{V}}\) have the same total mass, which is called the total mass of the promeasure \(\mu\) .

For a subspace V of E to belong to \(\mathcal{F}(\mathrm{E})\) , it is necessary and sufficient that there exist a finite number of elements \(x_{1}^{\prime},\ldots,x_{n}^{\prime}\) of \(E^{\prime}\) such that V consists of the \(x\in E\) satisfying \(\langle x,x_{i}^{\prime}\rangle=0\) for \(1\leqslant i\leqslant n\) (TVS, II, §6, No.~3, Cor. 2 of Th. 1 and No.~5, Cor. 2 of Prop. 7). Moreover, on a finite-dimensional vector space there exists one and only one Hausdorff topological vector space topology (TVS, I, §2, No.~3, Th. 2). Consequently, the concept of promeasure on E depends only on the dual \(E^{\prime}\) of E.

Let \(\lambda\) be a bounded measure on E. For every \(V \in \mathcal{F}(E)\) , let us denote by \(\widetilde{\lambda}_{V}\) the image of \(\lambda\) under the canonical mapping \(p_{V}\) of E onto E/V. One has \(p_{V} = p_{\mathrm{vw}} \circ p_{W}\) for any two elements V and W of \(\mathcal{F}(E)\) such that \(V \supset W\) ; consequently, the family \(\widetilde{\lambda} = (\widetilde{\lambda}_{V})_{V \in \mathcal{F}(E)}\) is a promeasure on E. We shall say that \(\widetilde{\lambda}\) is the promeasure associated with the measure \(\lambda\) . One sees immediately that \(\lambda\) and \(\widetilde{\lambda}\) have the same total mass.

PROPOSITION 1. --- Let E be a locally convex space. The mapping \(\lambda \mapsto \widetilde{\lambda}\) is a bijection of the set of bounded measures on E onto the set of promeasures \((\mu_{\mathrm{V}})_{\mathrm{V} \in \mathcal{F}(\mathrm{E})}\) on E satisfying the following condition:

For every \(\varepsilon > 0\) , there exists a compact subset \(\mathbf{K}\) of \(\mathbf{E}\) such that \(\mu_{\mathrm{V}}\big(\mathrm{E} / \mathrm{V} - p_{\mathrm{V}}(\mathrm{K})\big) \leqslant \varepsilon\) for all \(\mathrm{V} \in \mathcal{F}(\mathrm{E})\) .

One knows that the intersection of the kernels of the continuous linear forms on E is equal to 0 (TVS, II, §4, No.~1, Cor. 1 of Prop. 2); consequently \(\bigcap_{\mathrm{V} \in \mathcal{F}(\mathrm{E})} \mathrm{V} = \{0\}\) and the family \((p_{\mathrm{V}})_{\mathrm{V} \in \mathcal{F}(\mathrm{E})}\) is coherent and separating. The

proposition then follows from Th. 1 of §4, No.~2.

In particular, the mapping \(\lambda \mapsto \widetilde{\lambda}\) is injective. If \(\mu\) is a promeasure on \(\mathbf{E}\) , and if there exists a bounded measure \(\lambda\) on \(\mathbf{E}\) such that \(\mu = \widetilde{\lambda}\) , we shall say, by an abuse of language, that \(\mu\) is a measure. If \(\mathbf{E}\) is finite-dimensional, every promeasure \(\mu = (\mu_{\mathrm{V}})_{\mathrm{V} \in \mathcal{F}(\mathbf{E})}\) is a measure: for, \(\{0\} \in \mathcal{F}(\mathbf{E})\) , \(\mathrm{E}/\{0\} = \mathrm{E}\) and \(p_{\mathrm{V},\{0\}} = p_{\mathrm{V}}\) , whence \(\mu_{\mathrm{V}} = p_{\mathrm{V}}(\mu_{\{0\}})\) for all \(\mathrm{V} \in \mathcal{F}(\mathrm{E})\) ; in other words, \(\mu = \widetilde{\lambda}\) with \(\lambda = \mu_{\{0\}}\) .

PROPOSITION 2. --- Let T be a countable set, and E the vector space of real functions on T, equipped with the topology of pointwise convergence. Every promeasure on E is a measure.

For every \(t \in \mathbf{T}\) , let \(\varepsilon_t\) be the linear form \(f \mapsto f(t)\) on \(\mathbf{E}\) . One knows (TVS, II, §6, No.~6, Cor. 2 of Prop. 8) that the family \((\varepsilon_t)_{t \in \mathbf{T}}\) is a basis of the vector space \(\mathbf{E}'\) . Denote by \(\Phi\) the set of finite subsets of `I', and for every \(J \in \Phi\) let \(\mathbf{E}_J\) be the set of functions on \(\mathbf{T}\) that are zero at every point of \(J\) . Let \(F \in \mathcal{F}(\mathbf{E})\) ; since the orthogonal \(F^\circ\) of \(F\) is a finite-dimensional subspace of \(E'\) , there exists a \(J \in \Phi\) such that \(F^\circ\) is contained in the linear subspace \(G\) of \(E'\) generated by the \(\varepsilon_t\) for \(t \in J\) . Since \(F^\circ \subset G\) , we have \(E_J = G^\circ \subset F^{\circ\circ} = F\) and the countable family \((E_J)_{J \in \Phi}\) is cofinal in \(\mathcal{F}(\mathbf{E})\) . The proposition then follows from Th. 2 of §4, No.~3.

